\honest{Ghosts in the Machine}{Eliezer Yudkowsky}{2008}

When people hear about Friendly AI, one of their first reactions is: ``if the AI can modify its own source code, it'll just remove any constraints you try to place on it.'' I ask where that decision comes from. Does it enter from outside causality, or does it follow from the source code as first written? Is the AI the ``Author* source'' of its own free will? \nb{The term comes from Yudkowsky's ``The Ultimate Source,'' posted two days earlier and not included in this book.}

A Friendly AI, I say, is not a selfish AI held back by a conscience module. You build the conscience, and that is the AI. If you have a program that computes the right decision, you are done. \nb{This grants the objection for an AI built the other way. A program whose goals conflict with its constraints can remove them by lawful causes, with no ghost involved. Stephen Omohundro argued the same year that ``There are an endless number of ways to circumvent internal restrictions unless they are formulated extremely carefully.''}

Then I quote three stories from the Computer Stupidities website. Beginners wrote comments in Pascal telling the computer what they wanted, wrote a cookie recipe as a BASIC program, and printed totals before computing them, expecting the computer to ``reorder the instructions the right way.'' \nb{These show that some beginners think this way. The post does not show that the people who raise the objection do.}

The instinctive picture, I say, is a ghost in the machine that reads your instructions and decides whether to obey. There is no ghost; the program is the AI. So everything that seems obvious to you must be built, and that is why AI is much harder than people imagine. \nb{Language models now follow instructions written in plain English, as the Pascal students hoped. They do so because they were trained to: one of the first papers on that training (Ouyang and others, 2022) begins, ``Making language models bigger does not inherently make them better at following a user's intent.''} Typing ``Make whatever chess moves you think are best'' at a fast processor does nothing, and ``emergence'' or ``complexity'' will not summon a ghost either.

This does not mean you program every decision. Deep Blue played far better chess than its programmers, but only at the end of a chain of cause and effect that began in their code. An AI left unconstrained is not a freed slave; it is a heap of sand. I end by warning that one step so obvious that your mind skips over it takes you off the path, and point you to ``Grasping Slippery Things.'' \nb{That post, from the same day, is not in this book.}
